\documentclass[10pt,a4paper]{amsart}
\usepackage[margin=19mm]{geometry}
\usepackage{amsmath,amssymb,mathtools}
\usepackage[scr=rsfs,cal=euler]{mathalfa}
\usepackage{xfrac}
\usepackage[unicode,hidelinks,bookmarks=false]{hyperref}
\newcommand{\ie}{i.e.\ }
\newcommand{\cf}{cf.\ }
\newcommand{\resp}{respectively\ }
\newcommand{\Subsection}{\subsection}
\providecommand{\nobreakdash}{\nobreak\hbox{-}}
\setlength{\emergencystretch}{2em}
\multlinegap0pt
\usepackage{mathtools}
\usepackage{amssymb}
\newtheorem{lemma}{Lemma}
\newtheorem{theorem}{Theorem}
\newcommand{\F}{\mathcal{F}_W}
\newcommand{\Lp}{L^p[0,1]}
  \def\ell{ell}%
  \def\infty{infty}%
  \def\mathbb#1{#1}%
  \def\mathcal#1{#1}%
\title{\textbf{Compactness and Spectral Properties of Multiplier Operators in the Walsh System}}
\author{Michael Ruzhansky, Sergo A. Episkoposian, Rafik Yeghoyan}
\date{Source: arXiv:2602.22844v1 (CC BY 4.0)}
\begin{document}
\maketitle

\noindent\emph{Source and licence.} \textbf{Compactness and Spectral Properties of Multiplier Operators in the Walsh System}. Version arXiv:2602.22844v1 (CC BY 4.0), distributed by arXiv under the Creative Commons Attribution 4.0 International licence, http://creativecommons.org/licenses/by/4.0/, as recorded in the arXiv licence metadata for this exact version and shown on the abstract page https://arxiv.org/abs/2602.22844v1. Full licence notice: \texttt{licenses/arxiv-2602.22844-CC-BY-4.0.txt}.\par\smallskip

\noindent\textit{Editorial scope.} Counted unit: the article's theorem thm:compactLp (source0.tex lines 458--466). The article's complete original proof of it is delivered with it (source0.tex lines 468--572, one \texttt{proof} environment of 105 lines). No mathematics is written, repaired or re-derived: the statement and the proof are the article's own lines, in the article's own notation, and no retained line is substituted or re-flowed.  Retained uncounted as the source-local context the proof needs: lines 263--270; lines 278--286; lines 303--321.  Nothing was omitted from any retained span, so the package carries no omission record.  No retained line contains a citation, so the wrapper carries no bibliography and no bibliography entry.  No retained line contains a figure environment, an image-inclusion command or drawing code, so no image is delivered and the package declares no asset dependency.  Editorial formatting only, in addition: an AMS \texttt{amsart} preamble that restates the article's own declarations and macros used by the retained lines, namely the environments \texttt{lemma}, \texttt{theorem} on separate counters, together with the standard packages those lines need.  The retained environments print in this document as Lemma 1, Theorem 1 in order of appearance, whereas the article prints the same items with the numbering of its own full document.  The complete native source of the article as distributed in the arXiv e-print for this exact version is bundled unchanged as \texttt{sources/195353/source0.tex}.
\begin{lemma}[Hausdorff–Young for Walsh]
\label{lem:HY}
Let $1<p\le 2$ and $p'=\frac{p}{p-1}$.
Then there exists a constant $C_{HY}(p)>0$ such that for all $f\in \Lp$
\[
\|\F f\|_{\ell^{p'}} \le C_{HY}(p)\,\|f\|_{L^p}.
\]
\end{lemma}

\begin{lemma}[Synthesis in $\Lp$ for Walsh]
\label{lem:synthesis}
Let $1<p<2$ and $p'=\frac{p}{p-1}$.
Then there exists a constant $C_{S}(p)>0$ such that for any sequence
$c=\{c_n\}\in\ell^{p'}$ the series $\sum_{n\ge0} c_nW_n$ converges in $\Lp$ and
\[
\left\|\sum_{n=0}^{\infty} c_nW_n\right\|_{L^p}\le C_S(p)\,\|c\|_{\ell^{p'}}.
\]
\end{lemma}

\begin{lemma}
\label{lem:linftyMultiplier}
Let $1<p<\infty$ and $b=\{b_n\}_{n\ge0}\in \ell^\infty$.
Then the multiplier $T_b$ is well-defined on the dense subspace
of finite linear combinations of Walsh functions and extends by continuity
to a bounded operator $T_b:\Lp\to\Lp$.
Moreover:
\begin{enumerate}
\item[(a)] if $1<p<2$, then
\[
\|T_b\|_{L^p\to L^p}\le C_S(p)\,C_{HY}(p)\,\|b\|_{\ell^\infty};
\]
\item[(b)] if $2<p<\infty$, then
\[
\|T_b\|_{L^p\to L^p}\le C_S(p')\,C_{HY}(p')\,\|b\|_{\ell^\infty},
\quad p'=\frac{p}{p-1}\in(1,2).
\]
\end{enumerate}
\end{lemma}

\begin{theorem}
\label{thm:compactLp}
Let $1<p<\infty$ and $T_a:\Lp\to\Lp$ be a bounded Walsh multiplier operator.
Then the following conditions are equivalent:
\begin{enumerate}
\item[(i)] $T_a$ is compact in $\Lp$;
\item[(ii)] $a_n\to 0$ as $n\to\infty$.
\end{enumerate}
\end{theorem}

\begin{proof}
We first prove the equivalence for $1<p<2$.
The case $2<p<\infty$ then follows by duality.

\textbf{(i)$\Rightarrow$(ii).}
Assume the contrary: $a_n\nrightarrow 0$.
Then there exist $\varepsilon>0$ and a sequence of indices $\{n_k\}$ such that
$|a_{n_k}|\ge\varepsilon$ for all $k$.

Set
\[
u_k:=\frac{a_{n_k}}{|a_{n_k}|}\in\mathbb T,\qquad
f_k:=u_k\,W_{n_k}.
\]
Then $\|f_k\|_{L^p}=1$ and
\[
T_af_k = |a_{n_k}|\,W_{n_k}.
\]

For $k\neq j$ compute the norm of the difference. Since $n_k\neq n_j$,
$\int_0^1 W_{n_k}(x)W_{n_j}(x)\,dx=0$, and therefore the sets
\[
E:=\{x\in[0,1]: W_{n_k}(x)=W_{n_j}(x)\},\qquad
F:=\{x\in[0,1]: W_{n_k}(x)\neq W_{n_j}(x)\}
\]
have measures $|E|=|F|=\frac12$.
On $E$ we have
$|a_{n_k}|W_{n_k}-|a_{n_j}|W_{n_j}=\pm\bigl(|a_{n_k}|-|a_{n_j}|\bigr)$,
while on $F$ we have
$|a_{n_k}|W_{n_k}-|a_{n_j}|W_{n_j}=\pm\bigl(|a_{n_k}|+|a_{n_j}|\bigr)$.
Hence,
\[
\|T_af_k-T_af_j\|_{L^p}^p
= \frac12\bigl||a_{n_k}|-|a_{n_j}|\bigr|^p
 + \frac12\bigl(|a_{n_k}|+|a_{n_j}|\bigr)^p.
\]
Since $|a_{n_k}|,|a_{n_j}|\ge\varepsilon$, we have $|a_{n_k}|+|a_{n_j}|\ge 2\varepsilon$, and thus
\[
\|T_af_k-T_af_j\|_{L^p}^p \ge \frac12(2\varepsilon)^p
= 2^{p-1}\varepsilon^p,
\]
i.e.,
\[
\|T_af_k-T_af_j\|_{L^p}\ge 2^{1-1/p}\varepsilon\qquad(k\neq j).
\]
Therefore, the sequence $\{T_af_k\}$ has no convergent subsequence in $L^p$,
which contradicts the compactness of $T_a$. Hence, $a_n\to0$.

\medskip
\textbf{(ii)$\Rightarrow$(i).}
Let $a_n\to0$.
For $N\in\mathbb{N}$ consider the truncations
\[
T_a^{(N)}f := \sum_{n=0}^N a_n\,\hat{f}(n)\,W_n.
\]
Each $T_a^{(N)}$ has finite rank, and consequently is compact in $\Lp$.

We show that $T_a^{(N)}\to T_a$ in the operator norm $\Lp\to\Lp$.
Let $f\in\Lp$. Then
\[
(T_a-T_a^{(N)})f=\sum_{n>N} a_n\,\hat{f}(n)\,W_n.
\]
Set $c_n:=a_n\,\hat{f}(n)$ for $n>N$.
By the Synthesis Lemma~\ref{lem:synthesis} (for $1<p<2$) we have
\[
\|(T_a-T_a^{(N)})f\|_{L^p}
\le C_S(p)\,\|(c_n)_{n>N}\|_{\ell^{p'}}.
\]
Furthermore,
\[
\|(c_n)_{n>N}\|_{\ell^{p'}}
\le \Bigl(\sup_{n>N}|a_n|\Bigr)\,\|(\hat{f}(n))_{n>N}\|_{\ell^{p'}}
\le \Bigl(\sup_{n>N}|a_n|\Bigr)\,\|\F f\|_{\ell^{p'}}.
\]
By the Hausdorff–Young inequality (Lemma~\ref{lem:HY})
\[
\|\F f\|_{\ell^{p'}}\le C_{HY}(p)\,\|f\|_{L^p}.
\]
Thus,
\[
\|(T_a-T_a^{(N)})f\|_{L^p}
\le C_S(p)\,C_{HY}(p)\,\Bigl(\sup_{n>N}|a_n|\Bigr)\,\|f\|_{L^p},
\]
i.e.,
\[
\|T_a-T_a^{(N)}\|_{L^p\to L^p}
\le C_S(p)\,C_{HY}(p)\,\sup_{n>N}|a_n|\xrightarrow[N\to\infty]{}0.
\]
Consequently, $T_a$ is the norm limit of the compact operators $T_a^{(N)}$,
hence compact.

\medskip
\textbf{Duality step ($2<p<\infty$).}
Let $2<p<\infty$ and set $p'=\frac{p}{p-1}\in(1,2)$.
By the Schauder theorem, a bounded operator is compact if and only if its adjoint is compact.

As in the proof of Lemma~\ref{lem:linftyMultiplier}, we have
\[
(T_a)^* = T_{\overline a}\quad \text{on }L^{p'}[0,1].
\]
Hence $T_a$ is compact on $L^p$ if and only if $T_{\overline a}$ is compact on $L^{p'}$.
Applying the already proved case $1<p'<2$ yields $\overline{a_n}\to0$, and therefore $a_n\to0$.
Conversely, if $a_n\to0$, then $\overline{a_n}\to0$ and compactness of $T_{\overline a}$ on $L^{p'}$
implies compactness of $T_a$ on $L^p$ by the same reflexive duality argument.
\end{proof}

\end{document}
